Clearly, all pairs of the form $\paren{a, a}$ for $a \in \mathbb{N}$ are in $S$;
and for any $\paren{a, b} \in S$, also $\paren{b, a}$ is in $S$.

The pair $\paren{75, 1215}$ is in $S$, because
$75 = 3   \cdot 5^2$, $1215 = 3^5 \cdot 5$,
$76 = 2^2 \cdot 19$,  $1216 = 2^6 \cdot 19$.

The following is a family of elements of $S$,
parametrised by $n \in \mathbb{Z}^{+}$:
\begin{alignat*}{2}
 & a_n = 2^{n + 1} - 2 = 2 \cdot \paren{2^n - 1} \mspace{2mu} ,  && \qquad
 b_n = 2^{n + 1} a_n = 2^{2 n + 2} - 2^{n + 2}
\end{alignat*}
In fact,
$\operatorname{P}\paren{a} = \braces{2} \cup \operatorname{P}\paren{2^n - 1} \mspace{2mu}$
and $\mspace{2mu} \operatorname{P}\paren{b} = \braces{2} \cup \operatorname{P}\paren{a} = \operatorname{P}\paren{a}$;
and $b_n + 1 = \paren{2^{n + 1} - 1}^2 = \paren{a_n + 1}^2$,
so that $\operatorname{P}\paren{b + 1} = \operatorname{P}\paren{a + 1}$.

\vspace{2ex}
NOTE: This does not solve the problem yet.
